\begin{figure}[htbp]
  \centering
  \begin{adjustbox}{max width=\textwidth}
    \begin{tikzpicture}[
        % Gap ristretto a 0.8em
        node distance=1.2ex and 0.8em,
        %
        % Intestazione del canale
        ch-hdr/.style={draw=green!60!black, fill=green!18, rounded corners=1.5pt,
            minimum width=11.5em, minimum height=2.5em, font=\small\bfseries, align=center},
        %
        % Riga di metodo/dato dentro il canale
        ch-row/.style={draw=green!50!black, fill=green!7, rounded corners=1.5pt,
            minimum width=11.5em, minimum height=18pt,
            inner sep=4pt, font=\footnotesize\ttfamily, align=left},
        %
        % Blocco API
        ch-api/.style={draw=gray!40, fill=gray!10, rounded corners=1.5pt,
            minimum width=11.5em, minimum height=3.8em, inner sep=3pt, font=\footnotesize, align=left},
        %
        % Badge di fase
        ph-badge/.style={draw=gray!50, fill=gray!15, rounded corners=1.5pt,
            minimum width=11.5em, minimum height=3.5em, inner sep=3pt, font=\small, align=center},
      ]

      %% ─────────────────────────────────────────────────────────────────────
      %% Canale Token (sinistra)
      %% ─────────────────────────────────────────────────────────────────────
      \node[ch-hdr] (tok-hdr) {Canale Token};

      \node[ch-row, below=0.6ex of tok-hdr] (tok-r1) {\texttt{ROLE\_ADMIN}};
      \node[ch-row, below=0.6ex of tok-r1]  (tok-r2) {\texttt{ROLE\_USER\_A}};
      \node[ch-row, below=0.6ex of tok-r2]  (tok-r3) {\texttt{ROLE\_USER\_B}};

      \node[ch-api, below=0.6ex of tok-r3] (tok-api) {
        \texttt{set\_token(role, tok)}\\[1pt]
        \texttt{get\_token(role) \textrightarrow{} tok}\\[1pt]
        \texttt{has\_token(role) \textrightarrow{} bool}
      };

      \node[ph-badge, below=1.2ex of tok-api] (tok-ph) {\textbf{Fase~5:} test autenticati\\scrivono e leggono};

      %% ─────────────────────────────────────────────────────────────────────
      %% Canale Teardown LIFO (centro)
      %% ─────────────────────────────────────────────────────────────────────
      \node[ch-hdr, right=of tok-hdr] (td-hdr) {Canale Teardown \footnotesize[LIFO]};

      \node[ch-row, below=0.6ex of td-hdr] (td-r1) {push(fn\_C)};
      \node[ch-row, below=0.6ex of td-r1]  (td-r2) {push(fn\_B)};
      \node[ch-row, below=0.6ex of td-r2]  (td-r3) {push(fn\_A)};

      \node[ch-api, below=0.6ex of td-r3] (td-api) {
      \texttt{push\_teardown(callable)}\\[1pt]
      {\itshape drain order:}\\[1pt]
      \texttt{fn\_A \textrightarrow{} fn\_B \textrightarrow{} fn\_C}
      };

      \node[ph-badge, below=1.2ex of td-api] (td-ph) {\textbf{Fase~5:} push\\\textbf{Fase~6:} drain inv.};

      %% ─────────────────────────────────────────────────────────────────────
      %% Canale Shared Data (destra)
      %% ─────────────────────────────────────────────────────────────────────
      \node[ch-hdr, right=of td-hdr] (sh-hdr) {Canale Shared Data};

      \node[ch-row, below=0.6ex of sh-hdr] (sh-r1) {\texttt{"1.1.token"}};
      \node[ch-row, below=0.6ex of sh-r1]  (sh-r2) {\texttt{"1.1.repos"}};
      \node[ch-row, below=0.6ex of sh-r2]  (sh-r3) {\texttt{"ext.0.1.hits"}};

      \node[ch-api, below=0.6ex of sh-r3] (sh-api) {
        \texttt{set\_shared("{id}.{k}", v)}\\[1pt]
        \texttt{get\_shared("{id}.{k}")}\\[1pt]
        \texttt{~~\textrightarrow{} v}
      };

      \node[ph-badge, below=1.2ex of sh-api] (sh-ph) {\textbf{Fase~5:} produttore scrive\\batch successivi leggono};

      %% ── Sfondo per i tre canali ──────────────────────────────────────────
      \begin{scope}[on background layer]
        % Impostando inner sep=0.4em (esattamente la metà di 0.8em), 
        % i bordi delle colonne si toccano geometricamente senza sovrapporsi.
        \node[draw=green!50!black, fill=green!3, rounded corners=3pt, inner sep=0.4em,
          fit=(tok-hdr)(tok-r1)(tok-r2)(tok-r3)(tok-api)(tok-ph)] {};
        \node[draw=green!50!black, fill=green!3, rounded corners=3pt, inner sep=0.4em,
          fit=(td-hdr)(td-r1)(td-r2)(td-r3)(td-api)(td-ph)] {};
        \node[draw=green!50!black, fill=green!3, rounded corners=3pt, inner sep=0.4em,
          fit=(sh-hdr)(sh-r1)(sh-r2)(sh-r3)(sh-api)(sh-ph)] {};
      \end{scope}

      %% ── Intestazione globale (TestContext) ───────────────────────────────
      % La larghezza è esattamente 36.1em + 0.8em (margini esterni) = 36.9em.
      \node[draw=green!60!black, fill=green!25, rounded corners=2pt,
        minimum width=36.9em, font=\small\bfseries, align=center,
        above=3ex of td-hdr.north] (tc-master)
      {\texttt{TestContext} [mutable]: tre canali tipizzati con responsabilità distinte};

    \end{tikzpicture}
  \end{adjustbox}
  \caption{Struttura interna del \texttt{TestContext}: canali Token, Teardown \gls{LIFO} e Shared Data.}
  \label{fig:testcontext-channels}
\end{figure}